\label{kap:derivace}
\mathquote{It is unworthy of excellent men to lose hours like slaves in the labour of calculation which could safely be relegated to anyone else if machines were used.}{Gottfried Wilhelm Leibniz}

Derivace měří, jak rychle se něco mění. Rychlost je derivace dráhy podle času, zrychlení je derivace rychlosti. Sklon tečny ke grafu je derivace funkce. A když chceme najít maximum nebo minimum čehokoli, od zisku firmy po tvar plechovky, hledáme místa, kde je derivace nulová. Je to pravděpodobně nejužitečnější nástroj, který středoškolská matematika nabízí. Pokud vám z celé analýzy zůstane jen jedna věc, ať je to tahle.

\section{Pojem derivace}

Představte si auto, které ujede za čas $t$ dráhu $s(t)$. Průměrná rychlost mezi časy $t_0$ a $t_0 + h$ je
\begin{equation*}
  \frac{s(t_0 + h) - s(t_0)}{h}.
\end{equation*}
Okamžitou rychlost (tu, kterou ukazuje tachometr) dostaneme, když necháme $h$ jít k nule. Geometricky je průměrná rychlost směrnice \textbf{sečny} grafu procházející body $[t_0, s(t_0)]$ a $[t_0 + h, s(t_0 + h)]$. Když se druhý bod blíží k prvnímu, sečna se změní v \textbf{tečnu}.

\begin{center}
  \begin{tikzpicture}[scale=1]
    \draw[->] (-0.3,0) -- (4.6,0) node[right] {\small $x$};
    \draw[->] (0,-0.3) -- (0,4.3) node[above] {\small $y$};
    \draw[thick, colorofwisdom, domain=0.2:4.2, samples=60] plot (\x, {0.22*\x*\x + 0.3});
    \coordinate (A) at (1.5,0.795); \coordinate (B) at (3.5,2.995);
    \filldraw (A) circle (1.5pt); \filldraw (B) circle (1.5pt);
    \draw[gray] ($(A)!-0.4!(B)$) -- ($(A)!1.3!(B)$) node[right] {\small sečna};
    \draw[thick] (0.3,0.003) -- (4,2.445) node[right] {\small tečna};
    \draw[dashed] (A) -- (1.5,0) node[below] {\small $x_0$};
    \draw[dashed] (B) -- (3.5,0) node[below] {\small $x_0 + h$};
  \end{tikzpicture}
\end{center}

\begin{definitionbox}
  \textbf{Derivace funkce} $f$ v bodě $x_0$ je limita
  \begin{equation*}
    f'(x_0) = \lim_{h\to 0}\frac{f(x_0 + h) - f(x_0)}{h},
  \end{equation*}
  pokud existuje a je konečná. Pak říkáme, že $f$ je v bodě $x_0$ \textbf{diferencovatelná}. Funkce, která každému $x$ přiřadí $f'(x)$, je \textbf{derivace} $f'$.

  \vspace{0.2cm}

  \textbf{Geometrický význam}: $f'(x_0)$ je směrnice tečny ke grafu v bodě $[x_0, f(x_0)]$. \textbf{Rovnice tečny}:
  \begin{equation*}
    y - f(x_0) = f'(x_0)(x - x_0).
  \end{equation*}
  \textbf{Fyzikální význam}: derivace je okamžitá rychlost změny.
\end{definitionbox}

Leibniz derivaci značil $\frac{\mathrm{d}y}{\mathrm{d}x}$, tedy `nekonečně malá změna $y$ ku nekonečně malé změně $x$'. Newton psal tečku nad písmeno, $\dot{s}$. Značení $f'$ zavedl až Lagrange. Fyzici dodnes používají všechna tři, podle nálady.

\begin{example}
  Spočítejme derivaci $f(x) = x^2$ podle definice:
  \begin{equation*}
    f'(x) = \lim_{h\to 0}\frac{{(x + h)}^2 - x^2}{h} = \lim_{h\to 0}\frac{2xh + h^2}{h} = \lim_{h\to 0}(2x + h) = 2x.
  \end{equation*}
  Tečna ke grafu $y = x^2$ v bodě $x_0 = 3$ má směrnici $f'(3) = 6$ a rovnici $y - 9 = 6(x - 3)$, tedy $y = 6x - 9$.
\end{example}

\begin{note}
  Každá diferencovatelná funkce je spojitá. Obráceně to neplatí: funkce $|x|$ je v nule spojitá, ale derivaci tam nemá. Zprava je sklon $1$, zleva $-1$ a graf má `zlom'. Weierstrassova funkce z odbočky v kapitole~\ref{kap:funkce} je spojitá, ale nemá derivaci v žádném bodě. Je to funkce, která se skládá jen ze zlomů.
\end{note}

\section{Pravidla pro derivování}

\begin{definitionbox}
  \textbf{Derivace elementárních funkcí}:
  \begin{center}
    \renewcommand{\arraystretch}{2.3}
    \begin{tabular}{|c|c||c|c|}
      \hline
      $f(x)$ & $f'(x)$ & $f(x)$ & $f'(x)$ \\
      \hline
      $c$ & $0$ & $\sin x$ & $\cos x$ \\
      $x^n$ & $nx^{n-1}$ & $\cos x$ & $-\sin x$ \\
      $e^x$ & $e^x$ & $\tg x$ & $\dfrac{1}{\cos^2 x}$ \\
      $a^x$ & $a^x\ln a$ & $\cotg x$ & $-\dfrac{1}{\sin^2 x}$ \\
      $\ln x$ & $\dfrac{1}{x}$ & $\log_a x$ & $\dfrac{1}{x\ln a}$ \\
      \hline
    \end{tabular}
  \end{center}
  Vzorec pro $x^n$ platí pro libovolný reálný exponent $n$ (na vhodném definičním oboru), například $(\sqrt{x})' = {\left(x^{\frac{1}{2}}\right)}' = \frac{1}{2}x^{-\frac{1}{2}} = \frac{1}{2\sqrt{x}}$.

  \vspace{0.2cm}

  \textbf{Pravidla}:
  \begin{align*}
    (c\cdot f)' &= c\cdot f', & (f\pm g)' &= f'\pm g',\\
    (f\cdot g)' &= f'g + fg', & \left(\frac{f}{g}\right)' &= \frac{f'g - fg'}{g^2},\\
    \left(f(g(x))\right)' &= f'(g(x))\cdot g'(x) & &\text{(derivace složené funkce)}.
  \end{align*}
\end{definitionbox}

\begin{proof}[(několika vzorců)]
  \textbf{$(x^n)' = nx^{n-1}$ pro $n\in\mathbb{N}$}: Podle binomické věty z kapitoly~\ref{kap:kombinatorika} je ${(x + h)}^n = x^n + nx^{n-1}h + \binom{n}{2}x^{n-2}h^2 + \cdots$. Po odečtení $x^n$ a vydělení $h$ zbyde $nx^{n-1} + (\text{členy s } h)$, a ty pro $h\to 0$ zmizí.

  \textbf{$(\sin x)' = \cos x$}: Podle součtového vzorce
  \begin{equation*}
    \frac{\sin(x + h) - \sin x}{h} = \sin x\cdot\frac{\cos h - 1}{h} + \cos x\cdot\frac{\sin h}{h}.
  \end{equation*}
  Podle kapitoly~\ref{kap:limity} je $\lim\frac{\sin h}{h} = 1$. Dále $\frac{\cos h - 1}{h} = \frac{-2\sin^2\frac{h}{2}}{h} = -\sin\frac{h}{2}\cdot\frac{\sin\frac{h}{2}}{\frac{h}{2}}\to 0\cdot 1 = 0$. Celkem tedy $\cos x$. Tady je vidět, proč radiány: limita $\frac{\sin h}{h} = 1$ platí jen v nich. Ve stupních by derivace sinu byla $\frac{\pi}{180}\cos x$, což nechce nikdo.

  \textbf{Součin}: $f(x + h)g(x + h) - f(x)g(x) = \left(f(x + h) - f(x)\right)g(x + h) + f(x)\left(g(x + h) - g(x)\right)$. Po vydělení $h$ a limitě dostaneme $f'g + fg'$. (Použili jsme, že $g$ je spojitá, takže $g(x + h)\to g(x)$.)
\end{proof}

Teď dokážeme všechno ostatní z tabulky. Nejdřív zbylá pravidla, protože z nich pak vzorce pro jednotlivé funkce vyjdou skoro samy.

\begin{proof}[(zbylých pravidel)]
  \textbf{Konstanta, násobek a součet}: Pro $f(x) = c$ je $\frac{c - c}{h} = 0$, takže $c' = 0$. Pro $c\cdot f$ a $f\pm g$ stačí vytknout $c$, respektive rozdělit zlomek na dva, a použít, že limita násobku a součtu je násobek a součet limit (kapitola~\ref{kap:limity}).

  \textbf{Složená funkce}: Označme $k = g(x + h) - g(x)$. Pro $h\to 0$ jde i $k\to 0$, protože $g$ je spojitá. Pokud je $k\neq 0$, můžeme rozšířit:
  \begin{equation*}
    \frac{f(g(x + h)) - f(g(x))}{h} = \frac{f(g(x) + k) - f(g(x))}{k}\cdot\frac{g(x + h) - g(x)}{h}\ \longrightarrow\ f'(g(x))\cdot g'(x).
  \end{equation*}
  První zlomek je přesně definice derivace $f$ v bodě $g(x)$ (s přírůstkem $k$), druhý je definice $g'(x)$. (Pokud by $k$ bylo nulové pro libovolně malá $h$, musí se důkaz udělat trochu opatrněji, ale výsledek platí také. To si nechte na vysokou školu.)

  \textbf{Podíl}: Nejdřív $\left(\frac{1}{g}\right)'$. Funkce $\frac{1}{u} = u^{-1}$ má derivaci $-\frac{1}{u^2}$ (vzorec pro $x^n$ níže s $n = -1$), takže podle pravidla pro složenou funkci $\left(\frac{1}{g(x)}\right)' = -\frac{g'(x)}{g^2(x)}$. Pak podle pravidla pro součin
  \begin{equation*}
    {\left(\frac{f}{g}\right)}' = {\left(f\cdot\frac{1}{g}\right)}' = \frac{f'}{g} - \frac{fg'}{g^2} = \frac{f'g - fg'}{g^2}.
  \end{equation*}
\end{proof}

\begin{proof}[(zbylých vzorců z tabulky)]
  \textbf{$e^x$}: Podle definice
  \begin{equation*}
    \frac{e^{x + h} - e^x}{h} = e^x\cdot\frac{e^h - 1}{h},
  \end{equation*}
  takže stačí ukázat, že $\lim_{h\to 0}\frac{e^h - 1}{h} = 1$. Označme $t = e^h - 1$, tedy $h = \ln(1 + t)$, a pro $h\to 0$ jde i $t\to 0$. Pak
  \begin{equation*}
    \frac{e^h - 1}{h} = \frac{t}{\ln(1 + t)} = \frac{1}{\frac{1}{t}\ln(1 + t)} = \frac{1}{\ln{(1 + t)}^{\frac{1}{t}}}\ \longrightarrow\ \frac{1}{\ln e} = 1,
  \end{equation*}
  protože ${(1 + t)}^{\frac{1}{t}}\to e$, což je (pro $t = \frac{1}{n}$) přesně definice Eulerova čísla z kapitoly~\ref{kap:exponencialni}. Geometricky to znamená, že graf $e^x$ má v bodě $[0, 1]$ tečnu se směrnicí přesně $1$. Žádný jiný základ tuto vlastnost nemá, a proto je $e$ tak výjimečné.

  \textbf{$a^x$}: Protože $a = e^{\ln a}$, je $a^x = e^{x\ln a}$. Podle pravidla pro složenou funkci $(a^x)' = e^{x\ln a}\cdot\ln a = a^x\ln a$.

  \textbf{$\ln x$}: Pro $x > 0$ platí $e^{\ln x} = x$. Zderivujeme obě strany (vlevo podle pravidla pro složenou funkci): $e^{\ln x}\cdot(\ln x)' = 1$, tedy $x\cdot(\ln x)' = 1$ a $(\ln x)' = \frac{1}{x}$. Tomuto triku se říká \textbf{derivace inverzní funkce}: graf $\ln x$ je zrcadlový obraz grafu $e^x$ podle přímky $y = x$, takže i sklony tečen se `převrátí'.

  \textbf{$\log_a x$}: Podle převodu základu je $\log_a x = \frac{\ln x}{\ln a}$, a tedy $(\log_a x)' = \frac{1}{x\ln a}$.

  \textbf{$x^n$ pro libovolné reálné $n$} (pro $x > 0$): $x^n = e^{n\ln x}$, takže $(x^n)' = e^{n\ln x}\cdot\frac{n}{x} = x^n\cdot\frac{n}{x} = nx^{n-1}$.

  \textbf{$\cos x$}: Protože $\cos x = \sin\left(\frac{\pi}{2} - x\right)$, je $(\cos x)' = \cos\left(\frac{\pi}{2} - x\right)\cdot(-1) = -\sin x$.

  \textbf{$\tg x$ a $\cotg x$}: Podle pravidla pro podíl a vztahu $\sin^2 x + \cos^2 x = 1$:
  \begin{equation*}
    (\tg x)' = {\left(\frac{\sin x}{\cos x}\right)}' = \frac{\cos x\cos x - \sin x(-\sin x)}{\cos^2 x} = \frac{1}{\cos^2 x},\qquad(\cotg x)' = \frac{-\sin^2 x - \cos^2 x}{\sin^2 x} = -\frac{1}{\sin^2 x}.
  \end{equation*}
\end{proof}

Celou tabulku tedy stačí odvodit z definice derivace, ze dvou limit ($\frac{\sin h}{h}\to 1$ a definice $e$) a ze tří pravidel (součin, složená funkce a součet). Všechno ostatní je jen šikovné kombinování. Za zapamatování opravdu stojí jen pár řádků.

\begin{example}
  \begin{align*}
    (3x^4 - 5x^2 + 7)' &= 12x^3 - 10x,\\
    (x^2\sin x)' &= 2x\sin x + x^2\cos x,\\
    {\left(\frac{x + 1}{x - 1}\right)}' &= \frac{1\cdot(x - 1) - (x + 1)\cdot 1}{{(x - 1)}^2} = \frac{-2}{{(x - 1)}^2},\\
    \left(\sin(x^2)\right)' &= \cos(x^2)\cdot 2x,\\
    \left(e^{3x}\right)' &= e^{3x}\cdot 3,\\
    \left(\sqrt{x^2 + 1}\right)' &= \frac{1}{2\sqrt{x^2 + 1}}\cdot 2x = \frac{x}{\sqrt{x^2 + 1}},\\
    \left(\ln(\cos x)\right)' &= \frac{1}{\cos x}\cdot(-\sin x) = -\tg x.
  \end{align*}
  U složených funkcí derivujeme `zvenku dovnitř': nejdřív vnější funkci (vnitřek necháme, jak je) a pak vynásobíme derivací vnitřku. Rozklad na vnější a vnitřní funkci jsme si procvičili v kapitole~\ref{kap:funkce}.
\end{example}

\begin{problem}
  Zderivujte:
  \begin{enumerate}[label=(\alph*)]
    \item $f(x) = 2x^3 - \frac{1}{x} + \sqrt{x}$,
    \item $f(x) = xe^x$,
    \item $f(x) = \dfrac{x^2}{x + 1}$,
    \item $f(x) = {(2x - 1)}^5$,
    \item $f(x) = \ln(x^2 + 1)$,
    \item $f(x) = \sin^2 x$.
  \end{enumerate}
\end{problem}

\begin{solution}
  \begin{enumerate}[label=(\alph*)]
    \item $6x^2 + \frac{1}{x^2} + \frac{1}{2\sqrt{x}}$.
    \item $e^x + xe^x = (1 + x)e^x$.
    \item $\frac{2x(x + 1) - x^2}{{(x + 1)}^2} = \frac{x^2 + 2x}{{(x + 1)}^2}$.
    \item $5{(2x - 1)}^4\cdot 2 = 10{(2x - 1)}^4$. Roznásobovat pátou mocninu by byl hřích.
    \item $\frac{2x}{x^2 + 1}$.
    \item $2\sin x\cos x = \sin 2x$.
  \end{enumerate}
\end{solution}

\section{Průběh funkce}

\begin{definitionbox}
  \textbf{Monotonie}: Pokud je na intervalu $f'(x) > 0$, je $f$ na něm \textbf{rostoucí}. Pokud $f'(x) < 0$, je \textbf{klesající}.

  \textbf{Lokální extrémy}: Pokud má diferencovatelná funkce v bodě $x_0$ lokální extrém, je $f'(x_0) = 0$ (takovým bodům se říká \textbf{stacionární}). Pokud $f'$ v $x_0$ mění znaménko z $+$ na $-$, je tam \textbf{lokální maximum}, pokud z $-$ na $+$, \textbf{lokální minimum}. Ekvivalentně: pokud $f'(x_0) = 0$ a $f''(x_0) < 0$, je tam maximum, pokud $f''(x_0) > 0$, minimum.

  \textbf{Konvexnost a konkávnost}: Druhá derivace $f'' = (f')'$ udává, jak se mění sklon. Pokud $f'' > 0$, je funkce \textbf{konvexní} (graf `drží vodu', $\cup$). Pokud $f'' < 0$, je \textbf{konkávní} ($\cap$). Bod, ve kterém se konvexnost mění v konkávnost, je \textbf{inflexní bod}.
\end{definitionbox}

Pozor: $f'(x_0) = 0$ je jen \textit{nutná} podmínka, ne postačující. Funkce $x^3$ má $f'(0) = 0$, ale v nule extrém nemá, je všude rostoucí. Tečna je tam vodorovná, ale graf jde dál nahoru. Proto je třeba ověřit změnu znaménka derivace nebo druhou derivaci.

\begin{definitionbox}
  \textbf{Postup vyšetření průběhu funkce}:
  \begin{enumerate}
    \item definiční obor, sudost/lichost, periodicita,
    \item průsečíky s osami,
    \item limity v krajních bodech definičního oboru a asymptoty,
    \item první derivace: intervaly monotonie a lokální extrémy,
    \item druhá derivace: konvexnost, konkávnost a inflexní body,
    \item obor hodnot a graf.
  \end{enumerate}
\end{definitionbox}

\begin{example}
  Vyšetřeme $f(x) = x^3 - 3x$.
  \begin{enumerate}
    \item $D(f) = \mathbb{R}$. Funkce je lichá: $f(-x) = -x^3 + 3x = -f(x)$.
    \item $f(0) = 0$. Nulové body: $x(x^2 - 3) = 0$, tedy $x\in\{0, \pm\sqrt{3}\}$.
    \item $\lim_{x\to\infty} f(x) = \infty$, $\lim_{x\to-\infty} f(x) = -\infty$. Asymptoty nemá.
    \item $f'(x) = 3x^2 - 3 = 3(x - 1)(x + 1)$. Nulové body $\pm 1$. Na $(-\infty, -1)$ je $f' > 0$ (rostoucí), na $(-1, 1)$ je $f' < 0$ (klesající), na $(1, \infty)$ je $f' > 0$ (rostoucí). Lokální maximum $f(-1) = 2$, lokální minimum $f(1) = -2$.
    \item $f''(x) = 6x$. Konkávní na $(-\infty, 0)$, konvexní na $(0, \infty)$, inflexní bod $[0, 0]$.
    \item $H(f) = \mathbb{R}$.
  \end{enumerate}
  \begin{center}
    \begin{tikzpicture}[scale=0.8]
      \draw[->] (-3,0) -- (3.2,0) node[right] {\small $x$};
      \draw[->] (0,-3.2) -- (0,3.4) node[above] {\small $y$};
      \foreach \x in {-2,-1,1,2} \draw (\x,0.08) -- (\x,-0.08) node[below] {\tiny $\x$};
      \foreach \y in {-2,2} \draw (0.08,\y) -- (-0.08,\y) node[left] {\tiny $\y$};
      \draw[thick, colorofwisdom, domain=-2.2:2.2, samples=80] plot (\x, {\x*\x*\x - 3*\x});
      \filldraw (-1,2) circle (2pt) node[above] {\small max};
      \filldraw (1,-2) circle (2pt) node[below] {\small min};
      \filldraw (0,0) circle (2pt);
      \node[above right] at (0.05,0.05) {\tiny inflexe};
      \draw[dashed, gray] (-1.8,2) -- (-0.2,2);
      \draw[dashed, gray] (0.2,-2) -- (1.8,-2);
    \end{tikzpicture}
  \end{center}
\end{example}

\section{Užití derivací}

\subsection{Optimalizační úlohy}

V kapitole~\ref{kap:elementarni} jsme hledali maxima jen u kvadratických funkcí, kde stačilo najít vrchol. S derivacemi umíme najít extrém téměř čehokoli.

\begin{example}
  Plechovka tvaru válce má mít objem $1$ litr ($1000\ \si{cm^3}$). Jaké rozměry má mít, aby se na ni spotřebovalo co nejméně plechu?

  Z objemu $\pi r^2v = 1000$ máme $v = \frac{1000}{\pi r^2}$. Povrch je
  \begin{equation*}
    S(r) = 2\pi r^2 + 2\pi rv = 2\pi r^2 + \frac{2000}{r}.
  \end{equation*}
  Derivace $S'(r) = 4\pi r - \frac{2000}{r^2} = 0$, tedy $r^3 = \frac{500}{\pi}$ a $r\approx 5{,}42\ \si{cm}$. Výška vychází $v = \frac{1000}{\pi r^2}\approx 10{,}84\ \si{cm}$, tedy přesně $v = 2r$. Druhá derivace $S''(r) = 4\pi + \frac{4000}{r^3} > 0$, takže jde opravdu o minimum.

  Optimální plechovka má výšku rovnou průměru. Až půjdete příště do obchodu, všimněte si, že většina plechovek je vyšší. Na vině je prý marketing, pohodlný úchop a to, že víčka bývají z pevnějšího, a tedy dražšího plechu než plášť.
\end{example}

\begin{problem}
  \begin{enumerate}[label=(\alph*)]
    \item Z obdélníkového kartonu $30\times 20\ \si{cm}$ vystřihneme v rozích čtverce o straně $x$ a ohneme krabici bez víka. Pro jaké $x$ má největší objem?
    \item Najděte bod na parabole $y = x^2$, který je nejblíže bodu $[0, 2]$.
    \item Kámen vyhozený vzhůru je v čase $t$ ve výšce $h(t) = 20t - 5t^2$. Určete jeho rychlost v čase $t = 1$ a kdy dosáhne vrcholu.
  \end{enumerate}
\end{problem}

\begin{solution}
  \begin{enumerate}[label=(\alph*)]
    \item $V(x) = x(30 - 2x)(20 - 2x) = 4x^3 - 100x^2 + 600x$ pro $x\in(0, 10)$. $V'(x) = 12x^2 - 200x + 600 = 0$, tedy $3x^2 - 50x + 150 = 0$ a $x = \frac{50\pm\sqrt{2500 - 1800}}{6} = \frac{50\pm\sqrt{700}}{6}$. V intervalu je jen $x\approx 3{,}92\ \si{cm}$.
    \item Druhá mocnina vzdálenosti je $d^2 = x^2 + {(x^2 - 2)}^2 = x^4 - 3x^2 + 4$. Derivace $4x^3 - 6x = 2x(2x^2 - 3) = 0$. V bodě $0$ je lokální maximum ($d^2 = 4$), pro $x = \pm\sqrt{\frac{3}{2}}$ minimum s $d^2 = \frac{9}{4} - \frac{9}{2} + 4 = \frac{7}{4}$. Nejbližší body jsou $\left[\pm\sqrt{\frac{3}{2}}, \frac{3}{2}\right]$.
    \item $v(t) = h'(t) = 20 - 10t$, $v(1) = 10\ \si{m/s}$. Vrchol pro $v(t) = 0$, tedy $t = 2\ \si{s}$. Stejně jako v kapitole~\ref{kap:elementarni}, jen bez doplňování na čtverec. Zrychlení $h''(t) = -10\ \si{m/s^2}$ je tíhové zrychlení.
  \end{enumerate}
\end{solution}

\subsection{L'Hospitalovo pravidlo}

\begin{definitionbox}
  \textbf{L'Hospitalovo pravidlo}: Pokud je limita $\lim_{x\to a}\frac{f(x)}{g(x)}$ typu $\frac{0}{0}$ nebo $\frac{\infty}{\infty}$ a existuje $\lim_{x\to a}\frac{f'(x)}{g'(x)}$, pak
  \begin{equation*}
    \lim_{x\to a}\frac{f(x)}{g(x)} = \lim_{x\to a}\frac{f'(x)}{g'(x)}.
  \end{equation*}
\end{definitionbox}

\begin{example}
  \begin{equation*}
    \lim_{x\to 0}\frac{e^x - 1}{x} = \lim_{x\to 0}\frac{e^x}{1} = 1,\qquad\lim_{x\to\infty}\frac{\ln x}{x} = \lim_{x\to\infty}\frac{\frac{1}{x}}{1} = 0.
  \end{equation*}
  Druhá limita říká, že logaritmus roste pomaleji než jakákoli přímka. Pravidlo pochází od markýze de l'Hospital, který ho v roce 1696 publikoval v první učebnici diferenciálního počtu vůbec. Ve skutečnosti ho ale objevil jeho učitel Johann Bernoulli, kterému markýz platil za to, že mu bude posílat své objevy a nechá ho je vydávat pod svým jménem. Pravidlo tedy dodnes nese jméno člověka, který si ho koupil.
\end{example}

Na závěr malý výhled. Derivace odpovídá na otázku `znám funkci, jak rychle se mění?'. Často ale potřebujeme opak: `znám rychlost, kde jsem?' Na tuto obrácenou otázku odpovídá integrál. Že jde opravdu o obrácení derivace, a že tatáž operace zároveň počítá plochy pod grafem, říká \textbf{základní věta integrálního počtu} v sekci~\ref{sec:zvip} na straně~\pageref{sec:zvip}. Právě tato věta je to, co Newton a Leibniz objevili. Derivace samotné znali v nějaké podobě už Fermat a další.

\section{Příběh --- Válka o kalkulus}

Newton vymyslel diferenciální počet kolem roku 1666, během moru, kdy byla univerzita v Cambridge zavřená a on se uchýlil na rodinnou farmu. Nic ale nezveřejnil. Leibniz na totéž přišel nezávisle kolem roku 1675 a v roce 1684 to publikoval, s mnohem lepším značením, které používáme dodnes. Pak začala jedna z nejošklivějších hádek v dějinách vědy.

Newtonovi stoupenci obvinili Leibnize, že myšlenku ukradl. Leibniz se v roce 1711 obrátil na Královskou společnost v Londýně, aby spor rozhodla. Společnost jmenovala komisi, která po roce vydala zprávu, že prvenství patří Newtonovi a Leibniz je podezřelý z plagiátorství. Prezidentem Královské společnosti byl v té době Isaac Newton. Zprávu, jak se později ukázalo, napsal z velké části on sám. A pak o ní ještě anonymně napsal pochvalnou recenzi.

Leibniz zemřel v roce 1716 v Hannoveru, zapomenutý a v nemilosti u svého zaměstnavatele, který se mezitím stal anglickým králem Jiřím I. Na jeho pohřbu byl podle očitého svědka jediný člověk, jeho tajemník. Newton se o jeho smrti dozvěděl a prý prohlásil, že `zlomil Leibnizovi srdce'. Spor otrávil vztahy mezi britskými a kontinentálními matematiky na celé století. Britové z vlastenectví trvali na Newtonově nepraktickém značení s tečkami, a britská matematika proto za kontinentem desítky let zaostávala. Dnes víme, že oba přišli na kalkulus nezávisle, a celý svět používá Leibnizovo $\frac{\mathrm{d}y}{\mathrm{d}x}$.

\section{Shrnutí a otázky k ústní zkoušce}

\begin{problem}
  Odpovězte vlastními slovy:
  \begin{enumerate}
    \item Definujte derivaci funkce v bodě. Jaký má geometrický a fyzikální význam?
    \item Spočítejte derivaci $x^2$ podle definice. Jak vypadá rovnice tečny?
    \item Jaký je vztah mezi spojitostí a diferencovatelností? Uveďte příklad spojité funkce bez derivace.
    \item Uveďte derivace elementárních funkcí. Odvoďte derivaci $x^n$ a $\sin x$.
    \item Uveďte pravidla pro derivaci součtu, součinu, podílu a složené funkce.
    \item Jak souvisí znaménko první derivace s monotonií? Jak se hledají lokální extrémy? Je $f'(x_0) = 0$ postačující podmínka?
    \item Co říká druhá derivace? Co je inflexní bod?
    \item Popište postup vyšetření průběhu funkce.
    \item Uveďte příklad optimalizační úlohy a vyřešte ji.
    \item Co říká l'Hospitalovo pravidlo?
  \end{enumerate}
\end{problem}
